
\documentclass[../../main.tex]{subfiles}
\begin{document}

% 年度の大きな表示
\begin{center}
  {\fontsize{48pt}{58pt}\selectfont \textbf{1986}\normalsize\textbf{年度}}
\end{center}

\vspace{10mm}

% 本文

\vspace{15mm}

% 出題テーマと難易度の表
\noindent\textbf{出題テーマと難易度}

\vspace{5mm}

\begin{center}
  \shadowbox{%
    \begin{tabular}{|c|p{8cm}|c|}
      \hline
      \textbf{問題番号} & \textbf{テーマ} & \textbf{難易度} \\
      \hline
      第1問 & 整数問題（数列の公約数） & \\
      \hline
      第2問 & 空間図形と正四面体 & \\
      \hline
      第3問 & 微分方程式（関数決定） & \\
      \hline
      第4問 & 面積（放物線と3次曲線） & \\
      \hline
      第5問 & 面積と漸化式（ベータ関数） & \\
      \hline
    \end{tabular}%
  }
\end{center}

\vspace{15mm}

% 解答
\noindent\textbf{解答}

\vspace{5mm}

\begin{center}
  \shadowbox{%
    \renewcommand{\arraystretch}{1.6}
    \begin{tabular}{|c|c|p{8cm}|}
      \hline
      \textbf{問題番号} & \textbf{小問} & \textbf{答え} \\
      \hline
      第1問 & - & $7$ \\
      \hline
      \multirow{2}{*}{第2問} & (1) & $E\left(-\dfrac{1}{2},-\dfrac{1}{2},0\right)$，$F\left(\dfrac{1}{2},\dfrac{1}{2},1\right)$ \\
      \cline{2-3}
                             & (2) & 1辺の長さ $\sqrt{6}$ \\
      \hline
      第3問 & - & $f(x)=-x+2$ \\
      \hline
      第4問 & - & $\dfrac{64}{3}$ \\
      \hline
      \multirow{2}{*}{第5問} & (1) & $A(m,1)=\dfrac{1}{m+1}$ \\
      \cline{2-3}
                             & (3) & $A(m,n)=\dfrac{m!\,n!}{(m+n)!}$ \\
      \hline
    \end{tabular}%
    \renewcommand{\arraystretch}{1.0}
  }
\end{center}

\vspace{15mm}

% 解答時間
\noindent\textbf{試験形態}

\vspace{5mm}

\begin{center}
  \shadowbox{%
    \begin{tabular}{p{8cm}}
      （要確認）
    \end{tabular}%
  }
\end{center}

\end{document}
